% Chapter 9: Open Theoretical Problems and Mathematical Outlook
% Part I of the book.

\chapter{Open Theoretical Problems and Mathematical Outlook}
\label{ch:part1-open-theoretical-problems-and-mathematical}
\label{sec:conclusion}
% -----------------------------------------------------------------------------

The eight chapters before this one assemble a theory in three layers, and
each layer stops at a boundary the next cannot cross. \chapref{ch:part1-function-spaces-duality-and-sobolev-regula}
and \chapref{ch:part1-partial-differential-equations-as-solution} fixed the
stage: a partial differential equation (PDE) problem is a \emph{solution
operator} $\mathcal{G}:\mathcal{A}\to\mathcal{U}$ between infinite-dimensional
spaces, and the compact embeddings of Rellich--Kondrachov are what make images
of bounded sets small enough to
approximate at all.
\chapref{ch:part1-operator-approximation-theory-and-universa} and
\chapref{ch:part1-the-fourier-neural-operator-spectral-theor} supplied the
approximation guarantees --- Theorem~\ref{thm:chen_chen} in its classical form,
Theorem~\ref{thm:no_universality} for integral-kernel neural operators,
Theorem~\ref{thm:fno_convergence} for the Fourier neural operator (FNO),
\index{Fourier neural operator (FNO)!open problems}
whose error splits into a spectral truncation term of order $k_{\max}^{-(s+\alpha)}$
and a neural parametric term. \chapref{ch:part1-discretization-invariance-and-continuum-li}
then made exact the sense in which a trained network is a continuum object whose
grid realizations converge to it, \eqref{eq:limit_commute}, and
\chapref{ch:part1-physics-informed-operators-and-sobolev-reg} aligned the
training objective with those norms.

None of these results reaches the remaining terms in the error budget of an
actual surrogate. They are statements about \emph{approximation}, on a
\emph{compact set}, for \emph{one application}. They say nothing about how many
paired samples are needed to \emph{find} such a network from data; nothing about
what happens when the solution leaves the Sobolev scale the proofs require, as
at a shock; and nothing about what happens when the operator is applied not once
but $K = T/\Delta t$ times in feedback --- which is how every weather,
turbulence, or cardiac surrogate in this book is used. Those three
gaps are the subject of this chapter,\index{open problems} and they are not
matters of polish. The
first needs a learning theory in which parameter count no longer
tracks capacity and in which the samples are fields, not vectors; the
second needs a substitute for the precompactness that smoothing was providing,
because at a discontinuity the hypotheses of Theorem~\ref{thm:fno_convergence}
are violated rather than merely tight; the third needs dynamical systems results
--- absorbing sets, invariant measures, semicontinuity of attractors --- with no
counterpart in the approximation literature.

The reader should bring three pieces of background: empirical risk
minimization and covering numbers; entropy solutions of scalar conservation
laws; and the vocabulary of discrete dynamics (attractors, Lyapunov exponents,
invariant
measures). Section~\ref{sec:ch09-sample-complexity} treats sample complexity and
the ``curse of regularity'' that replaces the finite-dimensional curse of
dimensionality; Section~\ref{sec:ch09-shocks} computes how much of the
Part~I machinery survives a jump discontinuity; Section~\ref{sec:ch09-rollouts}
does the same for long autoregressive rollouts. This chapter introduces no new
architecture and proves no representation or convergence result that
Parts~II and~III depend on; its two labeled propositions are computations those
parts lean on and had not supplied. What the chapter does is state
what the book has not established, so that what follows is read with correct
expectations.

% -----------------------------------------------------------------------------
\section{Sample Complexity and the Curse of Regularity}
\label{sec:ch09-sample-complexity}
% -----------------------------------------------------------------------------

Fix an input probability measure $\mu$ on $\mathcal{A}$ and suppose we are handed
$n$ paired samples $\{(a_i,\mathcal{G}(a_i))\}_{i=1}^n$, each field stored on a
grid $\mathcal{M}_N$. Write $\mathcal{R}(\tilde{\mathcal{G}}) =
\mathbb{E}_{a\sim\mu}\|\tilde{\mathcal{G}}(a)-\mathcal{G}(a)\|^2_{L^2(\Omega)}$
for the population risk and $\hat{\mathcal{R}}$ for the empirical risk over
those $n$ samples. Fix a hypothesis class $\mathcal{H}$ of candidate operators,
and write $\mathcal{N}(\varepsilon,\mathcal{H},L^{2}(\mu))$ for the number of
balls of radius $\varepsilon$ in the $L^{2}(\mu)$ metric that suffice to cover
it. Symmetrization gives, up to an absolute constant and for
bounded or sub-Gaussian losses, the bound
\begin{equation}
  \|\hat{\mathcal{G}}-\mathcal{G}\|_{L^2(\mu)}
  \;\le\;
  2\underbrace{\inf_{\tilde{\mathcal{G}}\in\mathcal{H}}\|\tilde{\mathcal{G}}-\mathcal{G}\|_{L^2(\mu)}}_{\text{approximation}}
  \;+\;
  2\underbrace{\Big(\varepsilon+\sqrt{\tfrac{\log \mathcal{N}(\varepsilon,\,\mathcal{H},\,L^2(\mu))}{n}}\Big)}_{\text{estimation, any }\varepsilon>0}.
  \label{eq:ch09-risk}
\end{equation}
\index{sample complexity}\index{covering number}\index{estimation error!vs. approximation error}
Here $\hat{\mathcal{G}}$ is a minimizer of $\hat{\mathcal{R}}$ over
$\mathcal{H}$. The first term is precisely what
\chapref{ch:part1-operator-approximation-theory-and-universa}
and \chapref{ch:part1-the-fourier-neural-operator-spectral-theor} bound; nothing
more is needed there. The second term is where infinite dimensionality asserts
itself. In
$\mathbb{R}^d$ the $\varepsilon$-covering number of a ball scales as
$\varepsilon^{-d}$, so estimation error costs $\sqrt{d/n}$ up to logarithms: the
ambient dimension enters, and enters mildly. In a function space the analogous
quantity is finite only on sets that are compact in the metric being measured,
and for the unit ball $B_{H^s}$ of the Sobolev scale of
\chapref{ch:part1-function-spaces-duality-and-sobolev-regula}
the asymptotics are
\begin{equation}
  \log \mathcal{N}(\varepsilon, B_{H^s}, L^2) \sim \varepsilon^{-d/s}.
  \label{eq:ch09-entropy}
\end{equation}
\index{curse of dimensionality!in function space}\index{metric entropy}
Balancing the estimation and approximation contributions in
\eqref{eq:ch09-risk} --- setting
$\varepsilon \approx \varepsilon^{-d/(2s)}n^{-1/2}$ --- gives the accuracy a
uniform consistency argument can deliver from $n$ samples,
\begin{equation}
  \varepsilon_n \lesssim n^{-s/(2s+d)},
  \qquad\text{so}\qquad
  \mathcal{R}(\hat{\mathcal{G}}) = O\!\left(n^{-2s/(2s+d)}\right).
  \label{eq:ch09-rate}
\end{equation}
The relation is written one-sided on purpose: whether an estimator can do
better than \eqref{eq:ch09-rate} is exactly the second of the three gaps itemized
below, and two-sidedness
is not available. One silent premise also has to be named. The estimation term
in \eqref{eq:ch09-risk} is governed by the entropy of the \emph{hypothesis
class} $\mathcal{H}$ in $L^{2}(\mu)$, whereas the exponent $\varepsilon^{-d/s}$
inserted at the balance is the entropy \eqref{eq:ch09-entropy} of the
\emph{input} Sobolev ball; the derivation proceeds on the assumption that
$\mathcal{H}$ is chosen so that
$\log\mathcal{N}(\varepsilon,\mathcal{H},L^{2}(\mu)) \lesssim \varepsilon^{-d/s}$,
and no FNO or DeepONet parameterization is presently known to satisfy such a
condition --- the mismatch between what the risk bound controls and what the
entropy estimate computes is itself part of the open territory this chapter is
mapping.
The quantity that plays the role of the dimension is not $d$ but $d/s$: the
exponent is set by how many derivatives the data measure controls, not by how
many coordinates the PDE writes down. It is therefore better to speak
of a \emph{curse of regularity} than of a curse of dimensionality in function
space.\index{curse of regularity} In two dimensions, a class with one square
integrable derivative ($s=1$) yields $\varepsilon_n \lesssim n^{-1/4}$; a class with
jumps, which Section~\ref{sec:ch09-shocks} shows to lie in $H^s$ only for
$s<1/2$, yields $n^{-1/6}$; and as $s\downarrow 0$ the polynomial exponent
disappears.

\begin{assumption}[Compact data measure]
\label{ass:ch09-compact}
The measure $\mu$ in \eqref{eq:ch09-risk} is supported on a set $K\subset
\mathcal{A}$ whose $\varepsilon$-entropy in the input norm satisfies
\eqref{eq:ch09-entropy} for some $s>0$.
\end{assumption}
Assumption~\ref{ass:ch09-compact} is not a formality: it is the statistical twin
of the compactness hypothesis $K\subset\mathcal{A}$ compact in
Theorem~\ref{thm:no_universality}, and both are supplied by the same smoothing
properties catalogued in \chapref{ch:part1-partial-differential-equations-as-solution}.
For measures carrying only an $L^2$ moment ($s=0$) the right-hand side of
\eqref{eq:ch09-entropy} is infinite, no bound of the form \eqref{eq:ch09-rate}
is available, and one must fall back on non-uniform or Bayesian formulations, as
in \cite{stuart2010inverse}.

What exactly is missing comes in three pieces of different characters. \emph{First, upper bounds of the right shape.}
Approximation bounds expressed through the compactness of the input set exist for
both principal architectures \cite{kovachki2021universal,lanthaler2022error}, and
\cite{lanthaler2022error} also attaches a statistical term explicit in the
parameter count and in the Lipschitz constants of the realization. What is
lacking is a bound for the operators Parts~II and~III learn ---
viscous and inviscid Burgers, Darcy flow, two-dimensional Navier--Stokes
vorticity --- whose approximation term is the spectral rate of
Theorem~\ref{thm:fno_convergence} and whose estimation term is explicit in depth
$L$, channel count $d_v$, and mode cutoff $k_{\max}$ rather than in a generic
class $\mathcal{H}$ whose size the existence theorems leave unconstrained
(Remark~\ref{rmk:ch09-universality}). \emph{Second, lower bounds.} For finite-dimensional
nonparametric regression, matching minimax lower bounds are classical; for
operator classes nothing of the form ``no estimator given $n$ paired grid
samples from this class incurs error less than $\cdots$'' is known beyond one
exactly tractable model: the discrete-Fourier-transform least-squares estimator
of the linear FNO layer, for which \cite{subedi2024controlling} proves lower
bounds that match the truncation upper bound and fall strictly short of the
statistical and discretization upper bounds. For a nonlinear, multilayer
operator class no such statement is known, and without one we cannot tell
whether \eqref{eq:ch09-rate} is an artifact of the entropy
argument or a fact about the problem.\index{minimax lower bounds} \emph{Third, resolution-freeness.} The samples
entering the estimation term of \eqref{eq:ch09-risk} are stored at $N$ grid
points each, so a
naive reading suggests that refining the grid improves that term. But
the $N$ values of one realization are strongly dependent, and the number of
effectively independent quantities in $a|_{\mathcal{M}_N}$ is set by
\eqref{eq:ch09-entropy}, not by $N$. A bound in which the estimation term
depends on the parameter count and the entropy of $\mu$ but not on $N$ would
make the observed flatness of the error across mesh densities --- the identical
$0.0141$ at $N=64$ and $N=256$ quoted from Appendix~\ref{app:b} for a model trained
at a single resolution\index{discretization invariance!and generalization bounds} --- a
consequence rather than an observation.
There is structural reason to expect such a bound: the FNO parameter count
\begin{equation}
  P \asymp L\,d_v^{2}\,(2k_{\max}+1)^{d}
  \label{eq:ch09-paramcount}
\end{equation}
\index{parameter count!of a Fourier neural operator}
contains no grid size, as the implementation of
\chapref{ch:part2-fourier-neural-operators-fno-from-scratch} makes visible, and
the gap between the grid network and its continuum realization is separately
charged, at the algebraic rate \eqref{eq:ch06-rate} of
\chapref{ch:part1-discretization-invariance-and-continuum-li} --- a rate proved
there for the linear spectral layer and claimed, not proved, for the multi-layer
network. An architecture that attaches
a learnable bias per grid point carries $N$ degrees of freedom per
channel, so it forfeits the grid independence of \eqref{eq:ch09-paramcount} and
must be excluded from any such statement.

\begin{remark}[Universality without efficiency]
\label{rmk:ch09-universality}
Theorems~\ref{thm:chen_chen} and \ref{thm:no_universality} assert the
\emph{existence} of a network with error below $\varepsilon$; neither bounds the
$L$, $d_v$, or $k_{\max}$ needed to reach it, neither says the minimizer lies in
the region an optimizer searches, and neither constrains the estimation term of
\eqref{eq:ch09-risk} at all. Reading an existence theorem as a training
guarantee is the most common misuse of Part~I, and \eqref{eq:ch09-rate} is the
first place where a rate in $n$ rather than in $\varepsilon$ appears.
\index{universal approximation!limits of}
\end{remark}

% -----------------------------------------------------------------------------
\section{Discontinuities: Shock Propagation and the Gibbs Phenomenon}
\label{sec:ch09-shocks}
% -----------------------------------------------------------------------------

\chapref{ch:part1-partial-differential-equations-as-solution} observed that the
viscous Burgers equation $\partial_t u + u\,\partial_x u = \nu\,\partial_{xx}u$
on the torus $\mathbb{T}$ maps initial data $u_0\in L^2(\mathbb{T})$ into
$H^s(\mathbb{T})$ for every $s>0$ when $\nu>0$, and that as $\nu\to 0^+$ steep
gradients form and energy is transferred to arbitrarily high spatial
frequencies. With viscosity present, everything Part~I needs is available: the
operator smooths, the smoothing order $\alpha$ of
Theorem~\ref{thm:fno_convergence} is positive, and $k_{\max}^{-(s+\alpha)}$ decays.
In the inviscid limit the solution develops true
jump discontinuities and must be selected by an entropy condition, and the
machinery no longer applies;\index{entropy condition} the finite-volume literature is the reference for
this regime \cite{leveque2002finite}.

The failure is computable. Take the model discontinuity
$u = \mathbf{1}_{[0,1/2]}$ on $\mathbb{T}=[0,1)$ and let $P_{k_{\max}}$ denote
the $L^2$ orthogonal projection onto the modes $|k|\le k_{\max}$, the projection
used in the proof of Theorem~\ref{thm:fno_convergence}.

\begin{proposition}[Spectral tail of a jump]
\label{prop:ch09-jump}
For $u=\mathbf{1}_{[0,1/2]}$ one has $\hat{u}(0)=\tfrac12$ and, for $k\ne 0$,
$\hat{u}(k) = (1-e^{-\pi i k})/(2\pi i k)$,
hence $|\hat{u}(k)| = (\pi|k|)^{-1}$ for odd $k$ and $\hat{u}(k)=0$ for even
$k$. Consequently
\begin{equation}
  \|u - P_{k_{\max}} u\|_{L^2(\mathbb{T})}^{2}
  = \sum_{|k|>k_{\max}}|\hat{u}(k)|^{2}
  \;\le\; \frac{2}{\pi^{2}k_{\max}},
  \qquad\text{so}\qquad
  \|u-P_{k_{\max}}u\|_{L^2(\mathbb{T})} \le \frac{\sqrt{2}}{\pi}\,k_{\max}^{-1/2},
  \label{eq:ch09-jump-tail}
\end{equation}
and $u\in H^s(\mathbb{T})$ if and only if $s<\tfrac12$.
\end{proposition}
\begin{proof}
Direct integration over $[0,1/2]$ gives $\hat{u}(k)$, and
$|1-e^{-\pi i k}| = |1-(-1)^k|$ is $2$ for odd $k$ and $0$ for even $k$. Parseval
identifies the squared $L^2$ error with the tail sum, and
$\sum_{|k|>k_{\max}}k^{-2}\le 2/k_{\max}$ for $k_{\max}\ge1$ gives
\eqref{eq:ch09-jump-tail}.
Membership in $H^s$ requires
$\sum_k (1+k^2)^s|\hat u(k)|^2 \asymp \sum_{k\ \mathrm{odd}} k^{2s-2}<\infty$,
i.e. $2s-2<-1$.
\end{proof}
\index{Gibbs phenomenon}\index{spectral truncation!failure at discontinuities}

Two conclusions follow from Proposition~\ref{prop:ch09-jump}, of different
severity. The first is a rate:
$|\hat{u}(k)|\sim|k|^{-1}$ is the sharp decay a function of bounded variation can
be forced to have --- membership in $L^2$ alone asks only that
$\sum_k|\hat{u}(k)|^2$ converge, and fixes no pointwise rate at all --- and the
tail it leaves behind decays like $k_{\max}^{-1/2}$, so the exponent $s+\alpha$ of
Theorem~\ref{thm:fno_convergence} --- the output regularity of the target ---
must be read below $\tfrac12$ at a shock, and its constant $C_1$, priced by the
$H^{\sigma}$ norm of the target field for $\sigma<\tfrac12$, blows up as
$\sigma\uparrow\tfrac12$. The smoothing hypothesis itself fails outright: a
solution operator that ends in a jump gains no derivatives, so
\eqref{eq:ch05-error-decomposition} no longer applies and it is the tail
estimate that survives.
By \eqref{eq:ch09-entropy} the same penalty transfers to statistics: a data
measure whose samples contain fronts has small $s$, and the achievable
accuracy \eqref{eq:ch09-rate} degrades accordingly. The second conclusion is an
obstruction, not a rate. For $M\ge1$ let $S_M u=\mathcal{F}^{-1}(P_M\mathcal{F}u)$
denote the truncation of $u$ to the modes $|k|\le M$, the operator just used with
$M=k_{\max}$. Then for a piecewise smooth $u$ with jump
$[u]=u(x_0^+)-u(x_0^-)$ and midpoint $m=\tfrac12(u(x_0^+)+u(x_0^-))$,
\begin{equation}
  \lim_{M\to\infty}\ \sup_{x_0<x<x_0+c/M} S_M u(x)
  \;=\; m + [u]\cdot\frac{1}{\pi}\int_0^{\pi}\frac{\sin t}{t}\,dt
  \;=\; m+\Big(\tfrac12 + 0.08949\ldots\Big)[u],
  \label{eq:ch09-gibbs}
\end{equation}
for a suitable constant $c>0$; the Wilbraham--Gibbs overshoot\index{Wilbraham--Gibbs constant} is a fixed
$8.9\%$ of the jump height, confined to a layer of width $O(1/M)$ but of a
height that does not shrink with $M$. Since the
spectral branch of the FNO layer \eqref{eq:fno_layer} is exactly a trigonometric
polynomial of degree $k_{\max}$, it cannot converge in $L^\infty$ to a
discontinuous target. The pointwise bypass and the activation can carry a jump
forward from the input, so a trained network is not excluded from representing a
shock; what is excluded is doing so \emph{through the cheap, global,
$\mathcal{O}(N\log N)$ path} that makes the architecture efficient in the first
place. The Sobolev-regularized objective of
\chapref{ch:part1-physics-informed-operators-and-sobolev-reg} inherits it too:
$L^\infty$ is controlled only through an embedding requiring $s>d/2$,
the one regularity \eqref{eq:ch09-jump-tail} denies a discontinuous field.

Why, then, do the benchmarks of \chapref{ch:part2-canonical-pde-benchmarks-and-reference-wor}
report relative $L^2$ errors of order $1.5\%$ for Burgers --- $0.0139$ to $0.0160$, at
$\nu = 0.1$, on grids of $256$ to $8192$ points, for the solution at $t = 1$ alone
\cite{li2021fourier}? The honest answer is that that benchmark does not exercise the
failure just described. At $\nu = 0.1$ and $t = 1$ the heat factor
$e^{-\nu(2\pi m)^2 t}$ is $\approx e^{-3.95 m^2}$: the target the error is scored
against is a viscously smoothed field whose spectrum is negligible beyond a handful
of modes, with smooth inputs as well, so the flatness across resolutions is the
band-limited behavior of Proposition~\ref{prop:ch06-zeroshot}, not shock behavior
surviving in the metric. The metric-blindness argument belongs to the inviscid and
near-inviscid regime the cited columns do not test: an overshoot of fixed
height on a set of measure $O(1/M)$ contributes $O(M^{-1/2})$, vanishing as $M$
grows, and a unit jump displaced by $\delta$ contributes
$\|u(\cdot+\delta)-u\|_{L^2}=\sqrt{2\delta}$, small for sub-grid displacement.
Both defects are invisible in aggregate energy yet physically wrong --- a
saturation ringing outside $[0,1]$, an interface at the wrong place --- and no
$L^2$ number reports either; that is the regime of Part~III's finding
that a plain FNO struggles with sharp saturation fronts and needs a multi-scale
variant \cite{wen2022ufno}. \index{Burgers equation!inviscid limit}\index{shock wave}

The deepest difficulty is not the rate but the loss of compactness. The reason
\chapref{ch:part1-function-spaces-duality-and-sobolev-regula} can offer
precompact images is smoothing; the inviscid entropy solution operator supplies
no smoothing in $L^2$. What it does supply is a different pair of spaces: for
scalar conservation laws the entropy solution semigroup is a contraction in
$L^1$, and uniform bounds on total variation give precompactness in $L^1$ by
Helly's selection theorem; the finite-volume literature develops both facts
\cite{leveque2002finite}. The open problem is to
\index{bounded variation}
\emph{rebuild the precompactness hypothesis of
Theorem~\ref{thm:no_universality}, and with it the entropy estimate
\eqref{eq:ch09-entropy} that makes it quantitative, over the pair $(BV, L^1)$ instead of $(H^s,
L^2)$}, with three requirements attached, where $BV$ denotes the space of
functions of bounded variation. First, a compactness hypothesis that
is satisfied by the inviscid solution operators rather than their viscous
regularizations. Second, architectures allowed to output $BV$ functions without
ringing, which is what the three standard remedies aim at --- nonlinear
filtering, wavelet representations, adaptive framelet decompositions.
\index{wavelet representations}\index{framelet decompositions}
The asymmetry is worth stating exactly, because it is what an adaptive basis
buys: by \eqref{eq:ch09-jump-tail} a single jump reaches every band-limited mode,
whereas in a localized multiscale basis only the $O(1)$ basis functions whose
support meets the interface are affected at each scale, so of the $O(k_{\max})$
functions of a dictionary matched to the same finest resolution only
$O(\log k_{\max})$ are contaminated by the discontinuity. This is the regime
where adapting the representation is the lever. Adapting the mesh alone, as
graph-based operators do \cite{li2023geometry}, or reorganizing the
channel structure of a spectral layer by factorizing its multiplier array, buys
geometric or parameter-count flexibility but no additional regularity, and leaves
\eqref{eq:limit_commute} --- which asserts consistency and stays silent on
accuracy --- untouched. Third, a training signal that separates the entropy
solution from every other weak solution: the physics-informed residual loss \eqref{eq:pino_loss}
of \chapref{ch:part1-physics-informed-operators-and-sobolev-reg} penalizes
$\mathcal{N}(\mathcal{G}_\theta(a);a)$ in a negative or $L^2$-type norm, and
every weak solution satisfies that residual condition, so an entropy inequality
must be added; whether doing so buys convergence to the physically selected
solution is open.

% -----------------------------------------------------------------------------
\section{Infinite Time Horizons: Stability of Autoregressive Rollouts}
\label{sec:ch09-rollouts}
% -----------------------------------------------------------------------------

For evolutionary problems the operator of interest is the time-$\Delta t$ map
$\mathcal{G}_{\Delta t}: u(t)\mapsto u(t+\Delta t)$, applied recursively so that
$u(T) = \mathcal{G}_{\Delta t}^{K}u_0$ with $K=T/\Delta t$. Every application in
this book uses this form: the weather surrogate advances six-hour steps over
days, the cardiac model is unrolled over long temporal horizons to track
reentrant spiral waves, and the training
recipes of \chapref{ch:part2-training-regimes-optimization-and-failure-} exist
because naive feedback diverges. The pathwise error obeys an estimate that is
elementary and carries all the difficulty. Let $\mathcal{G}_\theta$ be the
trained network standing in for $\mathcal{G}_{\Delta t}$, let
$v_{k+1}=\mathcal{G}_\theta(v_k)$ be the network trajectory and $u(t_k)$ the exact
one at $t_k=k\Delta t$, let $e_k=\|v_k-u(t_k)\|$ with every norm below taken in
$L^2(\Omega)$, and let
\begin{equation}
  \tau \;=\; \sup_{k}\ \big\|\mathcal{G}_{\theta}\big(u(t_k)\big)-u(t_{k+1})\big\|
  \qquad\text{and}\qquad
  \Lambda_\theta \;=\; \sup_{v\ne w}\frac{\|\mathcal{G}_\theta(v)-\mathcal{G}_\theta(w)\|}{\|v-w\|}.
  \label{eq:ch09-onesided}
\end{equation}
The triangle inequality gives $e_{k+1}\le \Lambda_\theta e_k+\tau$, hence by
induction the discrete Gr\"onwall bound
\begin{equation}
  e_K \;\le\; \Lambda_\theta^{K}e_0 \;+\; \tau\,\frac{\Lambda_\theta^{K}-1}{\Lambda_\theta-1},
  \qquad\text{and}\qquad
  e_K \le e_0 + K\tau\ \ \text{if}\ \ \Lambda_\theta=1.
  \label{eq:ch09-gronwall}
\end{equation}
\index{autoregressive rollout}\index{Gr\"onwall inequality}
Everything useful, and nothing proved, lives in $\Lambda_\theta$: a supremum over
pairs of functions, as hard to evaluate as the dynamics it prices, whose upper
bound Proposition~\ref{prop:ch09-lip} makes an explicit functional of the
weights.

\begin{proposition}[An explicit Lipschitz bound for a band-limited realization]
\label{prop:ch09-lip}
Let $\mathcal{G}_\theta=\mathcal{Q}\circ\mathcal{V}_{L-1}\circ\cdots\circ\mathcal{V}_0\circ\mathcal{P}$
be an FNO $\mathcal{G}_\theta : L^2(\mathbb{T}^d;\mathbb{R}^{d_a}) \to
L^2(\mathbb{T}^d;\mathbb{R}^{d_u})$, with $\mathcal{V}_\ell$ the
layer of \eqref{eq:fno_layer}, elementwise activation of
Lipschitz constant $\ell_\sigma$, local maps $W_\ell$, spectral symbols
$R_{\phi,\ell}$ supported on $Z_{k_{\max}}$, and lift $\mathcal{P}$ and
projection $\mathcal{Q}$ with
Lipschitz constants $\Lambda_{\mathcal{P}}$ and $\Lambda_{\mathcal{Q}}$. Then, with
$\|\cdot\|_2$ the matrix spectral norm,
\begin{equation}
  \Lambda_\theta
  \;\le\;
  \Lambda_{\mathcal{P}}\Lambda_{\mathcal{Q}}
  \prod_{\ell=0}^{L-1}\ell_\sigma\Big(\|W_\ell\|_2+\sup_{k\in Z_{k_{\max}}}\|R_{\phi,\ell}(k)\|_2\Big).
  \label{eq:ch09-lipschitz}
\end{equation}
\end{proposition}
\begin{proof}
Plancherel turns the spectral branch into a fiberwise multiplication:
$\|\mathcal{F}^{-1}(R(k)\widehat{v}(k))\|_{L^2}^2=\sum_k\|R(k)\widehat{v}(k)\|_2^2
\le \sup_k\|R(k)\|_2^2\,\|v\|_{L^2}^2$. The bypass $W_\ell v$ contributes
$\|W_\ell\|_2\|v\|_{L^2}$, the bias is additive, and an elementwise Lipschitz
activation is Lipschitz on $L^2$ with the same constant. Multiplying the
per-layer constants and chaining the affine lift and projection proves
\eqref{eq:ch09-lipschitz}.
\end{proof}

\index{Lyapunov stability}\index{global attractor}
Equation~\eqref{eq:ch09-lipschitz} can be evaluated at any point in training, yet
no loss in this book controls it: weight decay penalizes the individual factors,
not the product across $L$ layers, and the objective is indifferent to
whether the composed map expands or contracts. It is tempting to answer by
requiring $\Lambda_\theta<1$, which by \eqref{eq:ch09-gronwall} bounds the error
uniformly in $K$. That answer is not merely conservative; it is
structurally wrong for the systems Part~III simulates.
\chapref{ch:part1-partial-differential-equations-as-solution} recorded, as a
hypothesis rather than a result proved there, that
two-dimensional Navier--Stokes flow in vorticity--streamfunction form possesses a
finite-dimensional global attractor\index{Navier--Stokes equations!global attractor} and maps bounded $L^2$ sets into compact
smooth subsets. Dissipativity of that kind contracts \emph{volumes} in phase
space, not \emph{distances}: the leading Lyapunov
exponents of the time-$\Delta t$ map are positive --- nearby states separate
exponentially, which is why forecasts have a horizon --- while their sum over
sufficiently many directions is negative, which is why trajectories settle onto
a finite-dimensional attractor. A network with $\Lambda_\theta<1$ cannot
reproduce the first half of that sentence. The requirement that the learned
discrete dynamical system inherit the dissipative structures and Lyapunov
stability of the original PDE attractor, without which climate and turbulence
forecasting is impossible, therefore cannot be met by a normwise contraction.
\index{dissipative estimate}

What is needed is dissipative rather than contractive, and the natural candidate
is an absorbing estimate in a norm the architecture can control:
\begin{equation}
  \big\|\mathcal{G}_\theta(v)\big\|_{H^r} \;\le\; a\,\|v\|_{H^r}+b,
  \qquad a<1
  \;\;\Longrightarrow\;\;
  \sup_{k\ge 0}\|v_k\|_{H^r}\le \max\Big(\|v_0\|_{H^r},\,\tfrac{b}{1-a}\Big).
  \label{eq:ch09-absorbing}
\end{equation}
The implication is a one-line iteration. Precompactness of the image, however, is
not free: the spectral branch is finite rank, but the local bypass $W_\ell v$
transmits a full-resolution copy of the state, so compactness has to be bought
from the estimate itself --- an absorbing bound in $H^{r}$ strictly above the
norm in which the dynamics are measured, then injected down by the
Rellich--Kondrachov compactness of
\chapref{ch:part1-function-spaces-duality-and-sobolev-regula}, is what makes the
image precompact. Given both, the classical criterion --- point dissipation plus
asymptotic compactness yields a compact global attractor --- gives the network
an attractor $\mathcal{A}_\theta$ inside the absorbing ball of
\eqref{eq:ch09-absorbing}. Two things
then remain unproved, and they are the problem. (i) The superposition map
$v\mapsto\sigma(v)$ is bounded on $H^r$ for
\index{superposition operator}
$0<r<1$ when $\sigma$ is Lipschitz, but at higher regularity composition demands
differentiability of $\sigma$ that the piecewise-linear activations of Part~II do
not have, though a smooth activation such as $\tanh$ does, so
\eqref{eq:ch09-absorbing} must be imposed in a norm strictly above the one in
which the dynamics are measured --- $L^2$ alone will not do, since an $L^2$
absorbing bound buys no precompactness --- and for the architectures of Part~II
the natural target is $H^1$. Whether the bound can be verified there is a
finite-dimensional computation about the composed network, for which
Proposition~\ref{prop:ch09-lip} already supplies the pieces; it is not the
unresolved regularity theory of three-dimensional Navier--Stokes, a question
about the PDE semigroup rather than about $\mathcal{G}_\theta$. (ii) Existence of
$\mathcal{A}_\theta$ says nothing about its relation to the PDE attractor
$\mathcal{A}$: upper semicontinuity, $\mathrm{dist}(\mathcal{A}_\theta,\mathcal{A})\to 0$,
requires the approximation to hold in a \emph{neighborhood} of $\mathcal{A}$ with
rate control on the approach to it, and Theorem~\ref{thm:no_universality}
constrains only the sup over a prescribed compact set for a single application;
lower semicontinuity --- that the learned
system does not lose part of the true attractor --- requires recurrence
properties of $\mathcal{G}_\theta$ that no architecture-level result provides.

\index{invariant measure}\index{statistical stability}
For a turbulent or climatic surrogate the reported quantity is typically not
$e_K$ at all but the long-time average of a scalar observable $\varphi$ of the
state, and the right error budget is then
\begin{equation}
  \left|\frac{1}{K}\sum_{k=0}^{K-1}\varphi(v_k)\;-\;\int\varphi\,d\nu\right|
  \;\le\;
  \underbrace{\left|\int\varphi\,d\nu_\theta-\int\varphi\,d\nu\right|}_{\text{bias of the learned invariant measure}}
  \;+\;
  \underbrace{\left|\frac1K\sum_k\varphi(v_k)-\int\varphi\,d\nu_\theta\right|}_{\text{fluctuation, governed by mixing}},
  \label{eq:ch09-weak}
\end{equation}
where $\nu$ is the physical invariant measure and $\nu_\theta$ one for the
learned map, furnished on the absorbing ball of \eqref{eq:ch09-absorbing} by a
Krylov--Bogolyubov argument, uniquely or not. The bias term in
\eqref{eq:ch09-weak} is a \emph{weak} error: it can be small when the pathwise
error \eqref{eq:ch09-gronwall} is exponentially large, because two chaotic
trajectories on the same attractor diverge while sharing its statistics. Making
this precise requires statistical stability of $\theta\mapsto\nu_\theta$,
established in detail for expanding one-dimensional maps and
essentially open for deterministic function-space maps that are neither uniformly
hyperbolic nor contractive.
Part~III's headline numbers sit on the wrong side of this gap: the
anomaly correlation reported for a six-hour step operator rolled out beyond day
seven \cite{pathak2022fourcastnet} is a spatial correlation of two fields at one
lead time, and nothing
in Part~I licenses extrapolating it to the $K\to\infty$ regime that climate
applications demand.\index{weather forecasting!error over long horizons}

The engineering response, developed in
\chapref{ch:part2-training-regimes-optimization-and-failure-} and
Appendix~\ref{app:d}, is
to train on the network's own unrolled trajectories rather than on teacher-forced
steps --- pushforward regularization,\index{pushforward regularization} which adds terms
$\sum_{k}\lambda_k\|\mathcal{G}_\theta^{k}(u_0)-u(k\Delta t)\|^2_{L^2}$ to the
loss. It works empirically, and the intuition is right: it moves the
guarantee from single applications of $\mathcal{G}_\theta$ to its
powers, which is where \eqref{eq:ch09-gronwall} says the error is born. What it
lacks is any theorem relating the weights $\lambda_k$, the unrolled horizon $K$,
the truncation error $\tau$ of \eqref{eq:ch09-onesided}, and the
Lipschitz product \eqref{eq:ch09-lipschitz}. Stating that relation --- what
condition on $(\lambda_k,K)$ makes training to horizon $K$ imply bounded error to
horizon $\alpha K$, or the existence of an invariant measure close to $\nu$ ---
is the most concrete of the three problems here and the one most tightly
constrained by available compute.

\takeaway{Summary of Part I}{Operator learning lifts machine learning from
Euclidean vectors to infinite-dimensional function spaces. Formalizing neural
operators as compositions of local linear maps and non-local integral kernels
buys approximation guarantees for the Fourier neural operator and DeepONet:
universality on compact sets, spectral convergence rates for smooth solution
operators, and discretization invariance in the continuum limit. The
corresponding statements about estimation from finite data, about discontinuous
solutions, and about infinite time horizons are the open problems catalogued
above.}

% -----------------------------------------------------------------------------
% Bridge to Part II
% -----------------------------------------------------------------------------
Part~I ended where the guarantees do: a network exists that approximates the
solution operator on a compact, smooth, single-step set, and we do not know how
many samples find it, what it does to a jump, or what it does after ten thousand
applications. Part~II of this book takes those three absences as its
specification. Opening with \chapref{ch:part2-the-practitioner-s-paradigm-infinite-dimen},
it turns the
continuum object into tensor programs that run on graphics processing units
(GPUs): batched channel-first layouts for fields, positional coordinate
embedding, relative $L^2$ losses that are invariant to grid density, and
pointwise normalization across an ensemble --- each the grid-level counterpart
of a continuum object from Part~I, and each a place where the
assumptions of this chapter become visible in code. The unrolled loss discussed
in Section~\ref{sec:ch09-rollouts} appears there as a training recipe rather than
a theorem, the truncation priced by \eqref{eq:ch09-jump-tail} appears as a choice
of \texttt{modes}, and the Lipschitz product \eqref{eq:ch09-lipschitz} appears as
a divergence the practitioner must notice by hand. Read Part~II, therefore, as
the empirical counterpart of a theory that is deliberately incomplete: what
follows is not a workaround for the gaps in \eqref{eq:ch09-risk},
\eqref{eq:ch09-gibbs}, and \eqref{eq:ch09-gronwall}, but the set of engineering
decisions one must make while those three problems remain open.

\vspace{1.5em}
